% !PRINTED 36
\weberpage{36}

Der erste Schritt ist die Bildung der Gleichung
\begin{equation}
 A={a_0\over b_0}B+C,
\tag{6}
\end{equation}
worin
\begin{equation}
 C=c_0x+c_1
\tag{7}
\end{equation}
und
\begin{equation}
 c_0={a_1b_0-b_1a_0\over b_0},\qquad
 c_1={a_2b_0-b_2a_0\over b_0}.
\tag{8}
\end{equation}
Ist $c_0=0$, also $C$ constant, so haben $A,B$ nur dann einen gemeinsamen Factor, wenn auch $c_1=0$ ist, und dann ist $A$ durch $B$ theilbar, d. h. $A$ und $B$ unterscheiden sich nur durch einen constanten Factor. Ist aber $c_0$ von Null verschieden, so setzen wir die Rechnung fort, indem wir
\[
 B=QC+D
\]
setzen, worin (nach \S~4, wenn dort $\alpha=-c_1:c_0$ gesetzt wird)
\begin{equation}
 D={b_0c_1^2-b_1c_0c_1+b_2c_0^2\over c_0^2}.
\tag{9}
\end{equation}
Ist $D$ von Null verschieden, so sind $A,B$ ohne gemeinschaftlichen Theiler. Ist aber $D=0$, so ist $C$ der grösste gemeinschaftliche Factor von $A$ und $B$.

Setzt man die Werthe $c_0,c_1$ aus (8) in (9) ein, so erhält man mit Weglassung des Nenners $b_0c_0^2$ die Bedingung für die Existenz eines gemeinsamen Theilers $C$ von $A$ und $B$ in der Form
\begin{equation}
\begin{aligned}
&a_0^2b_2^2+a_2^2b_0^2-2a_0a_2b_0b_2-a_1a_2b_0b_1-a_0a_1b_1b_2\\
&\qquad{}+a_0a_2b_1^2+a_1^2b_0b_2=0
\end{aligned}
\tag{10}
\end{equation}
oder
\begin{equation}
 (a_0b_2-b_0a_2)^2+(a_0b_1-a_1b_0)(a_2b_1-a_1b_2)=0.
\tag{11}
\end{equation}
Diese Bedingung ist auch erfüllt, wenn $c_0$ und $c_1$ gleich Null sind, und ist also die nothwendige und hinreichende Bedingung dafür, dass $A,B$ einen gemeinsamen Theiler haben.

Die linke Seite von (10) oder (11) heisst die \gesperrt{Resultante der Functionen $A$ und $B$}.

Als zweites Beispiel nehmen wir das schon im \S~3 gewählte. Setzen wir
\begin{equation}
\begin{aligned}
 f(x)&=a_0x^3+a_1x^2+a_2x+a_3=A,\\
 f'(x)&=3a_0x^2+2a_1x+a_2=B,
\end{aligned}
\tag{12}
\end{equation}
und setzen $a_0$ von Null verschieden voraus, so haben wir nach \S~3:
%CONT>>
